\subsection{\texorpdfstring{দশমিক থেকে সাধারণ $n$ ভিত্তি}{দশমিক থেকে সাধারণ n ভিত্তি}}
পূর্ণাংশে লক্ষ্য ভিত্তি দিয়ে ভাগ করো; ভাগশেষ নিচ থেকে ওপরে পড়ো। ভগ্নাংশে একই ভিত্তি দিয়ে গুণ করে পাওয়া অঙ্ক ওপর থেকে নিচে পড়ো।
\par\noindent\begin{minipage}{\linewidth}
\convExample{১}{পূর্ণসংখ্যা}
$ (37)_{10} $-কে ৩ ভিত্তিতে রূপান্তর করো।\par
\par\noindent\begin{minipage}{\linewidth}\centering
\phantomsection\label{VC068}
\begin{chThreeDiagram}\begin{tikzpicture}[x=1mm,y=1mm]
\path[use as bounding box] (-3,9) rectangle (72,-42);
\node[vcnote] at (15,7) {ভাজ্য/ভাগফল};
\node[vcnote] at (44,7) {ভাগশেষ};
\node[vcnum] at (0.0,0) {$ 3 $};
\node[vcnum] at (10.0,0) {$ 37 $};
\draw[vc line] (5.0,3)--(5.0,-3.5)--(50,-3.5);
\node[vcnum] (r0) at (44,-7) {$ 1 $};
\node[vcnum] at (1.5,-7) {$ 3 $};
\node[vcnum] at (11.5,-7) {$ 12 $};
\draw[vc line] (6.5,-4)--(6.5,-10.5)--(50,-10.5);
\node[vcnum] (r1) at (44,-14) {$ 0 $};
\node[vcnum] at (3.0,-14) {$ 3 $};
\node[vcnum] at (13.0,-14) {$ 4 $};
\draw[vc line] (8.0,-11)--(8.0,-17.5)--(50,-17.5);
\node[vcnum] (r2) at (44,-21) {$ 1 $};
\node[vcnum] at (4.5,-21) {$ 3 $};
\node[vcnum] at (14.5,-21) {$ 1 $};
\draw[vc line] (9.5,-18)--(9.5,-24.5)--(50,-24.5);
\node[vcnum] (r3) at (44,-28) {$ 1 $};
\node[vcnum] at (16.0,-28) {$0$};
\node[vcnote] at (61,-3) {LSB};
\node[vcnote] at (61,-33) {MSB};
\draw[vc arrow] (61,-28)--(61,-7);
\node[vcnote] at (31,-40) {নিচ থেকে ওপরে পড়ো};
\end{tikzpicture}\end{chThreeDiagram}
\end{minipage}\par
\[\therefore\ (37)_{10}=(1101)_{3}.\]
\end{minipage}\par
\par\noindent\begin{minipage}{\linewidth}
\convExample{২}{শুধু ভগ্নাংশ}
$ (0.875)_{10} $-কে ৪ ভিত্তিতে রূপান্তর করো।\par
\par\noindent\begin{minipage}{\linewidth}\centering
\phantomsection\label{VC069}
\begin{chThreeDiagram}\begin{tikzpicture}[x=1mm,y=1mm]
\path[use as bounding box] (-9,9) rectangle (76,-39);
\node[vcnote] at (14,7) {প্রাপ্ত অঙ্ক};
\node[vcnote] at (45,7) {ভগ্নাংশ};
\node[vcnum] at (45,0) {$ 0.875 $};
\draw[vc line] (26,3)--(26,-27);
\node[vcnum] (m0) at (45,-5) {$\times 4$};
\draw[vc line] (7,-8)--(65,-8);
\node[vcnum] at (14,-12) {$ 3 $};
\node[vcnum] (f0) at (45,-12) {$ 0.5 $};
\node[vcnum] (m1) at (45,-17) {$\times 4$};
\draw[vc line] (7,-20)--(65,-20);
\node[vcnum] at (14,-24) {$ 2 $};
\node[vcnum] (f1) at (45,-24) {$ 0 $};
\draw[vc arrow] (f0.east)--(72,-12)--(72,-17)--(m1.east);
\node[vcnote] at (0,-8) {MSB};
\node[vcnote] at (0,-29) {LSB};
\draw[vc arrow] (0,-12)--(0,-24);
\node[vcnote] at (34,-36) {ওপর থেকে নিচে পড়ো};
\end{tikzpicture}\end{chThreeDiagram}
\end{minipage}\par
\[\therefore\ (0.875)_{10}=(0.32)_{4}.\]
\end{minipage}\par
\par\noindent\begin{minipage}{\linewidth}
\convExample{৩}{মিশ্র সংখ্যা}
$ (42.6)_{10} $-কে ৫ ভিত্তিতে রূপান্তর করো।\par
\begin{minipage}[t]{.48\linewidth}\vspace{0pt}\centering\textbf{পূর্ণাংশ}\par\vspace{1mm}
\par\noindent\begin{minipage}{\linewidth}\centering
\phantomsection\label{VC070}
\begin{chThreeDiagram}\begin{tikzpicture}[x=1mm,y=1mm]
\path[use as bounding box] (-3,9) rectangle (72,-35);
\node[vcnote] at (15,7) {ভাজ্য/ভাগফল};
\node[vcnote] at (44,7) {ভাগশেষ};
\node[vcnum] at (0.0,0) {$ 5 $};
\node[vcnum] at (10.0,0) {$ 42 $};
\draw[vc line] (5.0,3)--(5.0,-3.5)--(50,-3.5);
\node[vcnum] (r0) at (44,-7) {$ 2 $};
\node[vcnum] at (1.5,-7) {$ 5 $};
\node[vcnum] at (11.5,-7) {$ 8 $};
\draw[vc line] (6.5,-4)--(6.5,-10.5)--(50,-10.5);
\node[vcnum] (r1) at (44,-14) {$ 3 $};
\node[vcnum] at (3.0,-14) {$ 5 $};
\node[vcnum] at (13.0,-14) {$ 1 $};
\draw[vc line] (8.0,-11)--(8.0,-17.5)--(50,-17.5);
\node[vcnum] (r2) at (44,-21) {$ 1 $};
\node[vcnum] at (14.5,-21) {$0$};
\node[vcnote] at (61,-3) {LSB};
\node[vcnote] at (61,-26) {MSB};
\draw[vc arrow] (61,-21)--(61,-7);
\node[vcnote] at (31,-33) {নিচ থেকে ওপরে পড়ো};
\end{tikzpicture}\end{chThreeDiagram}
\end{minipage}\par

$ (132)_{5} $
\end{minipage}\hfill
\begin{minipage}[t]{.48\linewidth}\vspace{0pt}\centering\textbf{ভগ্নাংশ}\par\vspace{1mm}
\par\noindent\begin{minipage}{\linewidth}\centering
\phantomsection\label{VC071}
\begin{chThreeDiagram}\begin{tikzpicture}[x=1mm,y=1mm]
\path[use as bounding box] (-9,9) rectangle (76,-27);
\node[vcnote] at (14,7) {প্রাপ্ত অঙ্ক};
\node[vcnote] at (45,7) {ভগ্নাংশ};
\node[vcnum] at (45,0) {$ 0.6 $};
\draw[vc line] (26,3)--(26,-15);
\node[vcnum] (m0) at (45,-5) {$\times 5$};
\draw[vc line] (7,-8)--(65,-8);
\node[vcnum] at (14,-12) {$ 3 $};
\node[vcnum] (f0) at (45,-12) {$ 0 $};
\node[vcnote] at (0,-10) {MSB=LSB};
\draw[vc arrow] (0,-14)--(0,-20);
\node[vcnote] at (34,-24) {ওপর থেকে নিচে পড়ো};
\end{tikzpicture}\end{chThreeDiagram}
\end{minipage}\par

$ (0.3)_{5} $
\end{minipage}\par
\[\therefore\ (42.6)_{10}=(132.3)_{5}.\]
\end{minipage}\par
